% =====================================================================
% Sistemas de Procesamiento para Big Data - UDP
% Presentacion de propuesta de proyecto (5 minutos / 5 diapositivas)
% Estructura: portada -> contexto -> problema -> relevancia -> datos
% Compilar con: pdflatex presentacion_fwi.tex  (dos veces)
% =====================================================================
\documentclass[11pt,aspectratio=169]{beamer}

\usepackage[utf8]{inputenc}
\usepackage[T1]{fontenc}
\usepackage[english]{babel}
\usepackage{booktabs}
\usepackage{amsmath}

% --- Tema y colores ---
\usetheme{Madrid}
\usecolortheme{whale}
\definecolor{primaryblue}{RGB}{24,76,120}
\setbeamercolor{structure}{fg=primaryblue}
\setbeamercolor{frametitle}{bg=primaryblue,fg=white}
\setbeamercolor{title}{bg=primaryblue,fg=white}
\setbeamertemplate{navigation symbols}{}
\setbeamertemplate{itemize items}[circle]

% --- Datos del documento ---
\title[Distributed FWI over Chile]{Distributed Computation of the Canadian Fire Weather Index (FWI)}
\subtitle{Decadal high-resolution gridded climate data in Chile}
\author[Group Members]{Agustín Pizarro \and Eduardo Valenzuela \and Benjamin Zuñiga}
\institute[UDP]{Big Data Processing Systems --- Universidad Diego Portales \\ \small Instructor: Nicolás Hidalgo}
\date{\today}

\begin{document}

% ----------------------- 1. Portada -----------------------
\begin{frame}[plain]
  \titlepage
\end{frame}

% ----------------------- 2. Contexto -----------------------
\begin{frame}{Context}
  \begin{itemize}
    \item Chile faces unprecedented wildfire seasons: \textbf{2017}, over 570,000 hectares burned;
          \textbf{2023}, over 430,000 hectares and dozens of casualties.
    \item Driven by 15+ years of megadrought, homogeneous flammable vegetation, and extreme days
          following the \textbf{30-30-30 rule} (T $>30^\circ$C, RH $<30\%$, wind $>30$ km/h).
    \item CONAF and SENAPRED operate on the \textbf{Canadian Fire Weather Index (FWI) System}:
          six interdependent sub-indices (FFMC, DMC, DC, ISI, BUI, FWI) that estimate fuel moisture
          and expected fire intensity from daily weather.
  \end{itemize}
\end{frame}

% ----------------------- 3. Problema -----------------------
\begin{frame}{The problem we want to address}
  \begin{itemize}
    \item Today the FWI is computed \textbf{only} at automatic weather stations. Chile's network is sparse,
          leaving \textbf{blind spots} in inner valleys and rural forest areas --- where recent megafires started.
    \item Covering the country with an interpolated $5$ km climate grid means computing
          $\sim$16,000 days $\times$ $\sim$20,000 cells $\approx \mathbf{3.2 \times 10^{8}}$ states.
    \item The model is \textbf{recursive}: day $t$ depends on day $t-1$, so it cannot simply be split by date.
      \[ FFMC_t = f(FFMC_{t-1},\, T_t,\, RH_t,\, W_t,\, P_t) \]
    \item A single machine fails: out-of-memory on 44-year matrices, fragmented NetCDF reads,
          and 30--50 hours per sequential run.
  \end{itemize}
\end{frame}

% ----------------------- 4. Relevancia -----------------------
\begin{frame}{Relevance}
  \begin{itemize}
    \item \textbf{Operational:} a nationwide daily fire danger map instead of a handful of isolated points,
          directly useful for prevention and resource allocation.
    \item \textbf{Scientific:} 44 years of reconstructed FWI enable trend analysis of fire danger
          under the megadrought.
    \item \textbf{Technical:} the case is a clean instance of a recurring Big Data challenge ---
          distributing \textbf{stateful, time-recursive models} over massive gridded data without
          saturating the network.
  \end{itemize}
  \vfill
  \begin{alertblock}{Research question}
    How to structure and partition a massive multi-decadal climate dataset on a distributed cluster
    (Apache Spark) to compute time-recursive models efficiently?
  \end{alertblock}
\end{frame}

% ----------------------- 5. Datos -----------------------
\begin{frame}{Data we plan to work with}
  \begin{table}
    \centering
    \small
    \begin{tabular}{ll}
      \toprule
      Provider   & CR2, Universidad de Chile (\texttt{vismet.cr2.cl}) \\
      Product    & \textbf{CR2MET v2.5} --- daily precipitation and max/min temperature \\
      Format     & NetCDF-4, arrays \texttt{[time, lat, lon]} \\
      Resolution & $0.05^\circ$ ($\approx 5$ km), continental Chile \\
      Coverage   & 1979--2023, 44 years of continuous daily records \\
      Volume     & Tens of GB compressed; $>320$ million states to compute \\
      Variety    & Raster grids + station tables for validation + land cover \\
      \bottomrule
    \end{tabular}
  \end{table}
  \vfill
  Processed with \textbf{Apache Spark} over \textbf{Apache Parquet}, partitioned by spatial cell so each
  executor holds the full time series it needs.
  \vfill
  \centering

\end{frame}

\end{document}